\section{TikZ}
\subsection{TikZとは}
TikZ(読み: ティクス，ティクツなど)とは，Beamerの作成者であるTill Tantau氏によって作成された描画パッケージである．
\subsection{TikZの導入方法}
TikZを使うときにLuaLaTeXの良さが現れる．LuaLaTeXでは他のパッケージと同様，以下のように宣言すれば良い．
\begin{lstlisting}
\usepackage{tikz}
\end{lstlisting}
一方，pLaTeXなどでは以下のようにパッケージの使用を宣言しなければならない．
\begin{lstlisting}
\usepackage[dvipdfmx]{graphicx, xcolor}
\usepackage{tikz}
\end{lstlisting}
これは，pLaTeXではtexファイルからDVIファイルという形式で書き出してから，PDFに変換するという都合上，dvipdfmxというドライバを指定する必要があるからだ．したがって，
グラフィック関係のパッケージを読み込むときにdvipdfmxドライバを指定するか，
documentclassを書くときにdvipdfmxドライバをオプションとして指定しなければならない．
\subsection{TikZ用のライブラリ}
TikZには専用のライブラリがあり，それらは\texttt{usetikzlibrary}コマンドで呼び出せる．多数あるため自由度が高いので，設定例を以下に示す．
\begin{lstlisting}
  \usepackage{tikz, tikz-3dplot, braket}
  \usetikzlibrary{patterns, intersections, calc, quotes, angles, arrows.meta, through, shapes.geometric}
  \usetikzlibrary{shapes, positioning}
\end{lstlisting}
\subsection{グラフの描き方}
TikZを使ってグラフを描くときは，\texttt{tikzpicture}環境を使う．
\begin{lstlisting}
\begin{tikzpicture}
  \draw (0,0) -- (1,1);
\end{tikzpicture}
\end{lstlisting}
このコードは，原点から(1,1)までの直線を描く．
\subsubsection{矢印の描き方}
矢印を描くときは，\texttt{->}または\texttt{-latex}を使う．
\begin{lstlisting}
\begin{figure}
  \begin{tikzpicture}
    \draw[->] (0,0) -- (1,1);
    \draw[-latex] (3,0) -- (4,1);
  \end{tikzpicture}
\end{figure}
\end{lstlisting}
\subsubsection{線の太さの変更}
線の太さを変更するときは，\texttt{line width}を使う．
\begin{lstlisting}
\begin{figure}
  \begin{tikzpicture}
    \draw[line width=2pt] (0,0) -- (1,1);
  \end{tikzpicture}
\end{figure}
\end{lstlisting}
\subsubsection{線の色の変更}
線の色を変更するときは，\texttt{color}を使う．
\begin{lstlisting}
\begin{figure}
  \begin{tikzpicture}
    \draw[color=red] (0,0) -- (1,1);
  \end{tikzpicture}
\end{figure}
\end{lstlisting}
\newpage
\subsubsection{線のスタイルの変更}
線のスタイルを変更するときは，\texttt{dashed}を使う．
\begin{lstlisting}
\begin{figure}
  \begin{tikzpicture}
    \draw[dashed] (0,0) -- (1,1);
  \end{tikzpicture}
\end{figure}
\end{lstlisting}
\subsubsection{線のパターンの変更}
線のパターンを変更するときは，\texttt{dotted}を使う．
\begin{lstlisting}
\begin{figure}
  \begin{tikzpicture}
    \draw[dotted] (0,0) -- (1,1);
  \end{tikzpicture}
\end{figure}
\end{lstlisting}
\subsubsection{線の角の変更}
線の角を変更するときは，\texttt{rounded corners}を使う．
\begin{lstlisting}
\begin{figure}
  \begin{tikzpicture}
    \draw[rounded corners] (0,0) -- (1,1);
  \end{tikzpicture}
\end{figure}
\end{lstlisting}
\subsubsection{線の始点と終点の変更}
線の始点と終点を変更するときは，\texttt{shift}を使う．
\begin{lstlisting}
\begin{figure}
  \begin{tikzpicture}
    \draw[shift={(1,1)}] (0,0) -- (1,1);
  \end{tikzpicture}
\end{figure}
\end{lstlisting}
\newpage
\subsubsection{線の透明度の変更}
線の透明度を変更するときは，\texttt{opacity}を使う．
\begin{lstlisting}
\begin{figure}
  \begin{tikzpicture}
    \draw[opacity=0.5] (0,0) -- (1,1);
  \end{tikzpicture}
\end{figure}
\end{lstlisting}
\subsubsection{点のプロット}
点をプロットするときは，\texttt{fill}を使う．
\begin{lstlisting}
\begin{figure}
  \begin{tikzpicture}
    \fill (0,0) circle [radius=1pt];
  \end{tikzpicture}
\end{figure}
\end{lstlisting}
\subsubsection{点のラベル}
点にラベルをつけるときは，\texttt{node}を使う．
\begin{lstlisting}
\begin{figure}
  \begin{tikzpicture}
    \node at (0, 1) [left] {A};
  \end{tikzpicture}
\end{figure}
\end{lstlisting}
\newpage
\subsubsection{まとめ}
以上のことをまとめると，以下のようなグラフを作成できる．
\begin{figure}[H]
  \centering
  \begin{tikzpicture}
    \draw[-latex] (-5,0) -- (5,0) node[right] {$x$};
    \draw[-latex] (0,-5) -- (0,5) node[above] {$y$};
    \draw[color=red, dashed, line width=2pt] (-5,-4) -- (4,5);
    \draw[color=blue, dotted, line width=2pt] (-5,3) -- (3,-5);
    \node at (0, 0) [above right] {O};
    \node at (0, 1) [left] {A};
    \node at (0, -2) [right] {B};
    \node at (-1.5, -0.5) [below] {C};
  \end{tikzpicture}
\end{figure}
\begin{lstlisting}
  \begin{figure}[H]
    \centering
    \begin{tikzpicture}
      \draw[-latex] (-5,0) -- (5,0) node[right] {$x$};
      \draw[-latex] (0,-5) -- (0,5) node[above] {$y$};
      \draw[color=red, dashed, line width=2pt] (-5,-4) -- (4,5);
      \draw[color=blue, dotted, line width=2pt] (-5,3) -- (3,-5);
      \node at (0, 0) [above right] {O};
      \node at (0, 1) [left] {A};
      \node at (0, -2) [right] {B};
      \node at (-1.5, -0.5) [below] {C};
    \end{tikzpicture}
  \end{figure}
\end{lstlisting}